\chapter{Testing and Evaluation}

Evaluating VOYO requires measuring two fundamentally different things: how well the
deterministic substrate performs (latency, correctness, edge-case robustness), and whether
tying an LLM to that substrate actually produces better plans than letting the LLM work alone.
Our evaluation strategy is structured around these two axes, drawing from AgentTravel's
published TravelBench framework which separates evaluation into knowledge-grounding metrics
and plan-quality metrics \cite{agenttravel_zhao_2024}. We map the former to our provenance,
retrieval, and latency tests, and the latter to our feasibility and A/B correctness tests.


\section{Unit Testing}

The core of our test suite consists of 99 deterministic unit tests across three subsystems, all
of which run without any live service dependency. The suite runs in 8.53 seconds and passes at
100\% on a clean Python 3.10 environment.

The \textbf{recommendation engine} is covered by 27 tests in
\texttt{tests/unit/recommendations/test\_engine.py}. These verify that the seven-dimensional
scoring function (category match, tag overlap, budget fit, pace fit, popularity, rating, and
recency boost) produces results proportional to the user profile provided. Our A/B correctness
scenarios confirm that a \texttt{history\_lover} profile surfaces historical POIs at the top of the
ranked list with scores between 0.56 and 0.64, while a \texttt{nature\_lover} profile surfaces
natural POIs scoring between 0.60 and 0.63, with no overlap in the top five results between the
two profiles. This constitutes measurable evidence that the engine is doing what it claims to do.

The \textbf{itinerary engine} is covered by 22 tests in
\texttt{tests/unit/itinerary/test\_itinerary\_engine.py}, verifying pace adjustment logic, cost
calculation, VROOM status parsing across all four recognised outcome codes
(\texttt{OPTIMAL}, \texttt{HEURISTIC}, \texttt{TIMEOUT}, and unknown), day-theme generation,
and the full pipeline integration. The \textbf{routing subsystem} is covered by 33 tests in
\texttt{tests/unit/routing/test\_routing.py}, verifying polyline decoding, Valhalla client
behaviour, opening-hours parsing across five different format variants (AM/PM, 24-hour,
``Open 24 hours'', \texttt{None}, and empty dict), POI-to-VROOM job conversion, VROOM problem
construction and solution parsing, and time string conversion.


\section{Edge Case and Robustness Testing}

A system that works on happy-path inputs but fails on malformed ones is not a system we can
defend at a presentation. We tested 18 explicit edge cases across all three deterministic
subsystems; all 18 pass in the 99-test core.

In the \textbf{recommendation engine}: an empty user profile degrades to defaults without
crashing; a POI with null fields (no tags, no price, no rating) is scored with neutral values on
missing dimensions rather than propagating a null error; an empty POI list returns an empty
recommendation set rather than an exception. In the \textbf{itinerary engine}: a POI with no
recorded visit duration defaults to 60 minutes; \texttt{None} ticket prices are skipped in cost
calculation without crashing; generating an itinerary from an empty POI list is handled
gracefully. In the \textbf{routing subsystem}: an empty polyline string returns an empty
coordinate list; any variant of malformed opening hours (\texttt{None}, \texttt{\{\}}, or missing
\texttt{weekday\_text} key) returns the default time window rather than an exception; an invalid
time string defaults to 09:00.

This kind of hardening is important specifically because VOYO's backend receives real-world
user input via CLEO and external database records — both of which are outside our control.
Every one of these edge cases represents a failure mode we observed or anticipated during
development and explicitly verified rather than assumed.


\section{Performance Testing}

Using synthetic data with no live service dependencies, we measured the latency for 13
benchmarks on the deterministic backend with three runs each, reporting the median and p95
across runs. The benchmark harness lives at \texttt{thesis/evidence/\_run\_benchmarks.py} and
was executed on 2026-06-15. Every benchmark passed its target threshold.

Scoring 200 POIs runs at a p95 of 1.662\,ms against a 200\,ms target — approximately 300
times under budget. VROOM problem construction runs at a median of 0.1016\,ms and a p95 of
0.1944\,ms against a 10\,ms target. VROOM solution parsing runs at a median of 0.0069\,ms
and a p95 of 0.0111\,ms. Opening-hours parsing runs at a median of 0.1048\,ms and a p95 of
0.2417\,ms against a 5\,ms target. Route polyline decoding runs at a median of 0.0053\,ms and
a p95 of 0.0123\,ms. CLEO context generation shows the most variability, with a median of
2.7807\,ms and a p95 of 45.972\,ms, though this comfortably passes its 100\,ms target.

These numbers substantiate what the methodology chapter argued: the deterministic substrate
is effectively free at our scale. The user's entire latency budget is consumed by CLEO (Groq
inference, typically 1--3 seconds per turn depending on query complexity) and live network
services such as Supabase and Valhalla. This is not a flaw — it is the intended consequence of
our architecture. We deliberately assigned all computationally heavy work to components that
scale independently of the LLM, because LLM latency is bounded by external infrastructure
that our code cannot control.


\section{Model Validation: CLEO Grounding Verification}

One of our most important engineering contributions is the force-tool grounding policy described
in §3.1.2. Claiming it works without testing it would be hollow. We ran 8 offline logic tests
covering every branch of the grounding policy without making any call to Groq.

Greeting inputs such as \texttt{hi} or \texttt{hello there} are correctly classified as
\texttt{concise} and receive no forced tool call. A general POI query such as ``Tell me about
Karnak Temple'' is classified as \texttt{standard} and forces any tool call, ruling out pure
memory answers. An itinerary query such as ``Plan a 2-day trip to Aswan'' is classified as
\texttt{detailed} and forces \texttt{search\_pois} specifically, guaranteeing database grounding
before any response is composed. POI name resolution (\texttt{\_resolve\_poi\_id}) correctly
converts a name string to a real database integer ID, passes an integer through unchanged,
and degrades gracefully on unrecognised inputs without crashing. All 8 cases pass.

We also ran a partial live verification before the Groq daily quota was exhausted. The same
itinerary query executed two genuine tool calls — \texttt{search\_pois} returning real database
records (id 171: Beit el-Wali, 100 EGP) and \texttt{curate\_itinerary} using the real returned IDs
— before the 429 quota error cut iteration 3 short. What matters is that these two real tool
calls prove the grounding path works end-to-end. Previously, the identical query returned a
hallucinated one-shot answer with zero tool calls. The grounding fix is not just asserted: it is
substantiated.


\section{Integration Testing}

We have integration test files for CLEO end-to-end behaviour
(\texttt{tests/integration/cleo/}), database interactions
(\texttt{tests/integration/database/}), and tool-cycle verification (\texttt{tests/tools/}). These
tests exist, are written, and are structured correctly — but their execution is gated on live
services and the Groq free-tier daily quota of 100,000 tokens. During test collection on
2026-06-15, the runner returned \texttt{groq.RateLimitError 429 — ``Limit 100000, Used
99855''}, confirming the quota constraint is the sole blocker, not a code defect.

One test file (\texttt{tests/integration/cleo/test\_cleo\_comprehensive.py}) additionally
contains a pre-existing syntax error in the test file itself (an unclosed f-string on line 24),
which is separate from the product code. We disclose this honestly: the 99-test core is the
complete verified surface of our system; the integration suite is the planned next step, gated
on upgrading to the Groq Developer Tier which removes the daily cap.


\section{Security Testing}

VOYO's security model rests on three layers. Supabase's built-in Row Level Security (RLS)
policies ensure that authenticated users can only query and modify their own data —
itineraries, conversation history, and profile records are isolated at the database level and
cannot be accessed across user accounts. JWT-based authentication validates every backend
request: the FastAPI middleware extracts and verifies the Supabase JWT before any route
handler executes, meaning unauthenticated requests are rejected before any database call is
made. The \texttt{service\_role} key used by the backend for administrative operations is never
exposed to the mobile client.

We did not conduct a formal penetration test or automated vulnerability scan, which we
acknowledge as a limitation appropriate to the scope of this thesis. What we can verify is that
the data isolation model is enforced at the database layer rather than only in application code
— meaning a bug in the API cannot circumvent it.


\section{Ablation Protocol (Pre-Registered — Measurement Pending)}

As registered in §3.5, our keystone experiment compares two configurations over identical
scenarios: \textbf{Configuration A} (full hybrid: CLEO + VROOM + Valhalla + geographic
coherence guard) and \textbf{Configuration B} (LLM-only: engines bypassed, candidate
filtering removed, LLM forced to author feasibility and routing). The three reported metrics are
feasibility \%, constraint-violation rate, and Average Margin — defined identically to Tang et al.\
\cite{itinera_tang_2024} to allow direct comparison against their published values of 86.0 for
the full system and 242.8 with the spatial optimiser removed.

We pre-register the pass threshold before the harness runs: hybrid feasibility $\geq$ 90\%,
LLM-only feasibility $\leq$ 50\%. This magnitude is derived from the ItiNera precedent
\cite[Table~2, p.~6]{itinera_tang_2024} and from TravelPlanner's independently quantified
GPT-4 success rate of 0.6\% on constrained planning \cite{travelplanner}. The evaluation
harness is built and the configurations are defined; measurement is pending stabilisation of the
live VROOM routing pipeline. We present the pre-registered protocol here because its value
lies in the commitment made before measurement, not in the result itself.
